\documentclass[10pt,a4paper]{amsart}
\usepackage[margin=19mm]{geometry}
\usepackage{amsmath,amssymb,mathtools}
\usepackage[scr=rsfs,cal=euler]{mathalfa}
\usepackage{xfrac}
\usepackage[unicode,hidelinks,bookmarks=false]{hyperref}
\newcommand{\ie}{i.e.\ }
\newcommand{\cf}{cf.\ }
\newcommand{\resp}{respectively\ }
\newcommand{\Subsection}{\subsection}
\providecommand{\nobreakdash}{\nobreak\hbox{-}}
\setlength{\emergencystretch}{2em}
\multlinegap0pt
\usepackage{bm}
\usepackage{amssymb}
\usepackage{enumerate}
\usepackage[shortlabels]{enumitem}
\newtheorem{lem}{Lemma}
\title{Bifurcation of limit cycles in a class of piecewise smooth generalized Abel equations with two asymmetric zones}
\author{Haihua Liang, Jianfeng Huang}
\date{Source: arXiv:2603.25149v1 (CC BY 4.0)}
\begin{document}
\maketitle

\noindent\emph{Source and licence.} Bifurcation of limit cycles in a class of piecewise smooth generalized Abel equations with two asymmetric zones. Version arXiv:2603.25149v1 (CC BY 4.0), distributed by arXiv under the Creative Commons Attribution 4.0 International licence, http://creativecommons.org/licenses/by/4.0/, as recorded in the arXiv licence metadata for this exact version and shown on the abstract page https://arxiv.org/abs/2603.25149v1. Full licence notice: \texttt{licenses/arxiv-2603.25149-CC-BY-4.0.txt}.\par\smallskip

\noindent\textit{Editorial scope.} Counted unit: the article's lem lemforSCD (source0.tex lines 408--432). The article's complete original proof of it is delivered with it (source0.tex lines 435--472, one \texttt{proof} environment of 38 lines). No mathematics is written, repaired or re-derived: the statement and the proof are the article's own lines, in the article's own notation, and no retained line is substituted or re-flowed.  Retained uncounted as the source-local context the proof needs: lines 399--404.  Nothing was omitted from any retained span, so the package carries no omission record.  No retained line contains a citation, so the wrapper carries no bibliography and no bibliography entry.  No retained line contains a figure environment, an image-inclusion command or drawing code, so no image is delivered and the package declares no asset dependency.  Editorial formatting only, in addition: an AMS \texttt{amsart} preamble that restates the article's own declarations and macros used by the retained lines, namely the environments \texttt{lem} on separate counters, together with the standard packages those lines need.  The retained environments print in this document as Lemma 1 in order of appearance, whereas the article prints the same items with the numbering of its own full document.  The complete native source of the article as distributed in the arXiv e-print for this exact version is bundled unchanged as \texttt{sources/197190/source0.tex}.
\begin{equation}\label{equa2-1}
\begin{split}
&\mathcal{C}_{k}^{E}(\rho,\beta):=\int_{E}\cos (k\theta)\left(1+\rho-\cos\theta\right)^{\beta}d\theta, \\ &\mathcal{S}_{k}^{E}(\rho,\beta):=\int_{E}\sin (k\theta)\left(1+\rho-\cos\theta\right)^{\beta}d\theta,
\\ &\mathcal{D}_{k}^{E}(\rho,\beta):=\int_{E}\theta\sin (k\theta)\left(1+\rho-\cos\theta\right)^{\beta}d\theta,
\end{split}
\end{equation}

\begin{lem} \label{lemforSCD} Let  $\mathcal{C}_{k}^{E}, \mathcal{S}_{k}^{E}$ and $\mathcal{D}_{k}^{E}$ be the functions defined by  \eqref{equa2-1}.
\begin{itemize}

\item[(i)] If  $\theta_1\in(0,\pi]$, then \begin{equation*}
\label{relaofscd-1}\mathcal{S}_{k}^{\left[\theta_{1},2\pi\right]}
=-\mathcal{S}_{k}^{\left[0,\theta_{1}\right]},\
\mathcal{C}_{k}^{\left[\theta_{1},2\pi\right]}
=\mathcal{C}_{k}^{\left[0,\theta_{1}\right]}+2\mathcal{C}_{k}^{\left[\theta_{1},\pi\right]},\
\mathcal{D}_{k}^{\left[\theta_{1},2\pi\right]}
=\mathcal{D}_{k}^{\left[0,\theta_{1}\right]}+2\mathcal{D}_{k}^{\left[\theta_{1},\pi\right]}-2\pi\left(\mathcal{S}_{k}^{\left[0,\theta_{1}\right]}+\mathcal{S}_{k}^{\left[\theta_{1},\pi\right]}\right).\end{equation*}

\item[(ii)]  If $\theta_1\in(\pi, 2\pi]$, then
\begin{equation*}\label{relaofscd-2}
 \mathcal{S}_{k}^{\left[0,\theta_{1}\right]}=\mathcal{S}_{k}^{\left[0,2\pi-\theta_{1}\right]}=-\mathcal{S}_{k}^{\left[\theta_{1},2\pi\right]},\quad
\mathcal{C}_{k}^{\left[0,\theta_{1}\right]}
=\mathcal{C}_{k}^{\left[0,2\pi-\theta_{1}\right]}+2\mathcal{C}_{k}^{\left[2\pi-\theta_{1},\pi\right]},\quad
\mathcal{C}_{k}^{\left[\theta_{1},2\pi\right]}
=\mathcal{C}_{k}^{\left[0,2\pi-\theta_1\right]},
\end{equation*}
\begin{equation*}\label{relaofscd-2-1}
 \mathcal{D}_{k}^{\left[\theta_{1},2\pi\right]}=\mathcal{D}_{k}^{\left[0, 2\pi-\theta_{1}\right]}-2\pi\mathcal{S}_{k}^{\left[0, 2\pi-\theta_{1}\right]},\quad
 \mathcal{D}_{k}^{\left[0, \theta_{1}\right]}=\mathcal{D}_{k}^{\left[0, 2\pi-\theta_{1}\right]}+2\mathcal{D}_{k}^{\left[ 2\pi-\theta_{1}, \pi\right]}-2\pi\mathcal{S}_{k}^{\left[ 2\pi-\theta_{1}, \pi\right]}.
\end{equation*}
\end{itemize}
\end{lem}

\begin{proof} (i)
We only provide the proof of the third equality since the first and the second equalities can be dealt with in a similar way. By the transformation $\theta \to 2\pi-\theta$, we get
\begin{equation*} \label{relaprooeq1}
\begin{split}
\mathcal{D}_{k}^{\left[\theta_{1},2\pi-\theta_1\right]}=&\mathcal{D}_{k}^{\left[\theta_{1},\pi\right]}+\mathcal{D}_{k}^{\left[\pi,2\pi-\theta_1\right]}\\
=&\mathcal{D}_{k}^{\left[\theta_{1},\pi\right]}+\int_{\theta_1}^{\pi}(\theta-2\pi)\sin(k\theta)\left(1+\rho-\cos\theta\right)^{\beta}d\theta\\
=&2\mathcal{D}_{k}^{\left[\theta_{1},\pi\right]}-2\pi\mathcal{S}_{k}^{\left[\theta_{1},\pi\right]},
\end{split}
\end{equation*}
and
\begin{equation*}
\mathcal{D}_{k}^{\left[2\pi-\theta_{1},2\pi\right]}=\int_0^{\theta_1}(\theta-2\pi)\sin(k\theta)\left(1+\rho-\cos\theta\right)^{\beta}d\theta=\mathcal{D}_{k}^{\left[0,\theta_{1}\right]}-2\pi\mathcal{S}_{k}^{\left[0,\theta_{1}\right]}.
\end{equation*}
Thus, $$\mathcal{D}_{k}^{\left[\theta_{1},2\pi\right]}=\mathcal{D}_{k}^{\left[\theta_{1},2\pi-\theta_1\right]}+\mathcal{D}_{k}^{\left[2\pi-\theta_{1},2\pi\right]}
=\mathcal{D}_{k}^{\left[0,\theta_{1}\right]}+2\mathcal{D}_{k}^{\left[\theta_{1},\pi\right]}-2\pi\mathcal{S}_{k}^{\left[0,\theta_{1}\right]}-2\pi\mathcal{S}_{k}^{\left[\theta_{1},\pi\right]}.$$



 (ii)
We only provide the proof of the   fourth and fifth  equalities. Using  the transformation $\theta \to 2\pi-\theta$, we get
\begin{equation*} \label{relaprooeq2}
\mathcal{D}_{k}^{\left[0,\theta_{1}\right]}=\int_{2\pi-\theta_1}^{2\pi}(\theta-2\pi)\sin(k\theta)\left(1+\rho-\cos\theta\right)^{\beta}d\theta=\mathcal{D}_{k}^{\left[0,2\pi\right]}-\mathcal{D}_{k}^{\left[0,2\pi-\theta_{1}\right]}+2\pi\mathcal{S}_{k}^{\left[0,2\pi-\theta_{1}\right]},
\end{equation*}
thus \begin{equation*}
\mathcal{D}_{k}^{\left[\theta_{1}, 2\pi\right]}=\mathcal{D}_{k}^{\left[0,2\pi\right]}-\mathcal{D}_{k}^{\left[0,\theta_{1}\right]}=\mathcal{D}_{k}^{\left[0,2\pi-\theta_{1}\right]}-2\pi\mathcal{S}_{k}^{\left[0,2\pi-\theta_{1}\right]}.
\end{equation*}

Next we prove the fifth equality. Again, by using the transformation $\theta \to 2\pi-\theta$,
\begin{equation*}
\mathcal{D}_{k}^{\left[\pi,\theta_{1}\right]}=\int_{2\pi-\theta_1}^{\pi}(\theta-2\pi)\sin(k\theta)\left(1+\rho-\cos\theta\right)^{\beta}d\theta=\mathcal{D}_{k}^{\left[2\pi-\theta_{1},\pi\right]}-2\pi\mathcal{S}_{k}^{\left[2\pi-\theta_{1},\pi\right]}.
\end{equation*}
Hence,
 \begin{equation*}
\mathcal{D}_{k}^{\left[0,\theta_{1}\right]}=\mathcal{D}_{k}^{\left[0,2\pi-\theta_1\right]}+\mathcal{D}_{k}^{\left[2\pi-\theta_1,\pi\right]}+\mathcal{D}_{k}^{\left[\pi,\theta_{1}\right]}=\mathcal{D}_{k}^{\left[0,2\pi-\theta_1\right]}+2\mathcal{D}_{k}^{\left[2\pi-\theta_{1},\pi\right]}-2\pi\mathcal{S}_{k}^{\left[2\pi-\theta_{1},\pi\right]}.
\end{equation*}


The proof is complete. \end{proof}

\end{document}
